%% kapa: exam=1 round=27 year=2016 period=1 form=A subject=real_estate subject_ko=부동산학원론 seq=6 number=86
%% source: Q-Net 감정평가사 기출문제 (cst00309) artlSeq=5169402 file=extracted/1008/5169402/제27회 감정평가사 1차 문제/제27회 감정평가사 1교시문제(A형).pdf
%% conversion: kordoc 4.24.51 -> tools/gen1.py (auto); review=proofread-sonnet; flags=none
\begin{question}{86}
부동산 중개계약에 관한 설명으로 옳은 것을 모두 고른 것은?
\begin{framedtext}
ㄱ. 독점중개계약: 매각의뢰를 받은 경우 그 계약기간 내에 거래가 성사되면 개업공인중개사가 당해 부동산거래를 성사시키지 않았더라도 중개수수료 청구권이 발생한다. \par
ㄴ. 전속중개계약: 공인중개사법령상 중개의뢰인은 중개대상물의 중개를 의뢰함에 있어서 특정한 개업공인중개사를 정하여 그 개업공인중개사에 한하여 당해 중개대상물을 중개하도록 하는 계약을 체결하여야 한다고 규정하고 있다. \par
ㄷ. 일반중개계약: 소유자는 다수의 개업공인중개사에게 매도를 의뢰할 수 있고, 매수인과의 거래를 먼저 성사시킨 개업공인중개사에게 수수료를 지불한다. \par
ㄹ. 공동중개계약: 부동산정보센터나 부동산협회 등을 매체로 하여 다수의 개업공인중개사가 상호 협동하여 공동으로 중개 역할을 하는 것을 말한다. \par
ㅁ. 순가중개계약: 거래가격을 정하고 이를 초과한 금액으로 거래가 이루어진 경우 초과액은 개업공인중개사와 의뢰인이 나누어 갖는 것이다. \par
\end{framedtext}
\begin{choices}
  \item ㄷ
  \item ㄱ, ㄴ
  \item ㄷ, ㄹ
  \item ㄱ, ㄷ, ㄹ
  \item ㄷ, ㄹ, ㅁ
\end{choices}
\end{question}
\answer{4}
